\begin{frame}{실용: 몇 차원을 남길 것인가 (스크리 플롯)}
  \begin{itemize}
    \item \kb{누적} 설명 분산 비율을 그려, ``더 남겨도 거의 의미 없는
      작은 고유값''부터는 언제인지를 고르는 것이 표준.
    \item 이 누적 곡선을 \kb{스크리 플롯}(scree plot,
      ``벼랑''이라는 뜻 --- 곡선이 급격히 떨어졌다가 평평해짐)이라
      부른다.
    \[ \text{누적 설명 분산} = \frac{\sum_{i=1}^{k}\lambda_{i}}
       {\sum_{i=1}^{d}\lambda_{i}} \]
    \item 가장 흔한 기준: \textbf{누적이 90\%(또는 95\%)를 넘는 가장
      작은 $k$}.
    \item 붓꽃의 누적: $[92.5\%,\; 97.8\%,\; 99.5\%,\; 100\%]$ ---
      1개만 남겨도 92.5\%, 2개면 97.8\%. ``95\% 기준''이면 $k = 2$.
  \end{itemize}
  \vskip0.6em
  {\footnotesize 같은 실용적 기준: 7.2절 랜덤 포레스트에서 트리 개수를
  ``더 늘려도 성능이 거의 안 오르는 지점''으로 정하던 것.}
\end{frame}
